\documentclass[11pt]{article}
\usepackage[margin=1in]{geometry}
\usepackage{parskip}
\usepackage{xcolor}
\usepackage{titlesec}

\titleformat{\section}[block]{\large\bfseries\color{blue!50!black}}{}{0pt}{}
\setlength{\parskip}{0.6em}

\title{Speaker Script: Extended-MFGAN talk\\
       \large 30-minute version, simple English}
\author{Hikmatullo Ismatov}
\date{}

\begin{document}
\maketitle

\section*{[Slide 1 -- Title page]}

Hello everyone. Thank you for being here today.

My name is Hikmatullo Ismatov. I am from the King Abdullah University of Science
and Technology, KAUST. This is joint work with my supervisor, Professor Diogo
Gomes.

The title of our paper is \emph{Generative Adversarial Networks for Mean Field
Games with Weakly Non-Separable Hamiltonians}. It is a long title, but I will
explain every word in the next 30 minutes.

The main message of the talk is this: there are two big families of deep
learning methods for Mean Field Games, and until now they could not be
combined. We show how to combine their strengths, and we prove a very clean
convergence rate.

\section*{[Slide 2 -- Roadmap]}

Let me start by telling you the structure of the talk. There are three parts.

In the first part, I will explain the problem. Why are deep learning methods
for Mean Field Games stuck in two camps that cannot talk to each other? One
camp is GAN-based, the other is physics-informed. Each one has a serious
weakness.

In the second part, I will present our solution. It is a contraction theorem
in the Wasserstein-2 metric, with an explicit and sharp convergence rate.
The rate is given by a simple formula: epsilon times $L_R$, divided by lambda.
I will explain what these symbols mean.

In the third part, I will show three numerical experiments. In all three,
the empirical convergence rate matches our theory to four decimal places.

The talk is 30 minutes total. Please feel free to ask questions at any time.

\section*{[Slide 3 -- What is a Mean Field Game?]}

So, what is a Mean Field Game?

The idea is the following. Imagine a very large population of agents. Each
agent is rational, which means each agent wants to minimise some cost. But
the cost depends on what the rest of the population is doing. So the agents
interact through the population distribution.

When we look for a Nash equilibrium, two things happen at the same time.

First, each agent solves an optimisation problem. This gives a value function,
called phi. The value function tells you the cost from now until the end of
the game. Phi solves a Hamilton-Jacobi-Bellman equation, the HJB equation,
and it goes backward in time.

Second, the population evolves. The density of agents is called rho. Rho
solves a Fokker-Planck equation, and it goes forward in time.

The two equations are coupled. The gradient of phi tells the agents how to
move. So phi sets the drift in the FP equation. And the density rho enters
the running cost in the HJB. So rho changes the optimisation.

The function $H$ is the Hamiltonian. It captures the cost and the interactions.
The parameter nu is the diffusion. Larger nu means more random noise.

Let me give you a concrete example, so you can keep it in mind. Imagine
many cars on a road. Each driver wants to reach a destination quickly, but
also wants to avoid traffic. The state $x$ is the position on the road,
the density rho is how crowded the road is, and the value function phi
tells each driver the best strategy. The Mean Field Game gives you the
equilibrium where every driver chooses optimally, given what all the other
drivers are doing.

This system has many applications: autonomous traffic, financial markets,
energy grids, opinion dynamics, and more recently multi-agent reinforcement
learning.

The mathematical theory of Mean Field Games is by now mature. The challenge
is computational. Solving the coupled HJB and Fokker-Planck system in high
dimensions is very hard. Classical finite-difference methods scale badly
with dimension. This is where deep learning comes in.

\section*{[Slide 4 -- Two families of deep MFG solvers]}

Now, deep learning methods for Mean Field Games come in two families.

On the left side, we have GAN-based methods. The most famous is APAC-Net,
and another one is MFGAN. These methods learn the density rho through a
\emph{generator} network, like a Wasserstein GAN. They give you the density
in a distributional sense. You can sample from it, you have a Wasserstein
convergence guarantee. This is very strong.

But these methods have one big weakness. They only work when the Hamiltonian
is \emph{separable}. Separable means that $H$ splits as $H_0$ of $x$ and $p$,
\emph{minus} $f_0$ of $x$ and rho. The dependence on momentum and density
do not mix. This is not just a convenience. It is a hard structural condition.
The mathematical formula they use, called the Hopf formula, breaks if the
Hamiltonian is not separable.

On the right side, we have physics-informed methods. The most famous is MFDGM,
which is the deep Galerkin method applied to Mean Field Games. These methods
just minimise the residual of the PDE, using a neural network. They handle
\emph{any} Hamiltonian, separable or not. So they are very general.

But they pay a price. They give you only $L^p$ convergence with $p$ less than
two, and they give no explicit rate. You also lose the distributional
representation of the density.

So we have a dichotomy. Power, or generality, but never both. If you want
strong convergence guarantees, you must restrict yourself to separable
Hamiltonians. If you want generality, you must give up the rate.

This is not a happy situation. Many real applications are non-separable.
And without an explicit rate, you do not know how many iterations to run,
or how to compare two methods. So this is a real gap in the literature.

Our paper resolves this dichotomy in a precise regime that we call
\emph{weakly non-separable}.

\section*{[Slide 5 -- The hard example: MFG-LWR]}

Let me show you the canonical hard example. It is the traffic flow model
of Huang, sometimes called MFG-LWR.

The Hamiltonian is $H$ of $x$, rho, $p$ equals one half $p$ squared, minus
the term $1$ minus rho times $p$.

Why is this important? It models traffic on a highway. The density rho is
how many cars are at position $x$. The momentum $p$ is related to their speed.
The term $1$ minus rho times $p$ couples them. When the road is full, the
cars slow down. So density and momentum interact in a non-trivial way.

Because of this term, you cannot apply APAC-Net. The Hopf formula breaks.
You can apply MFDGM, but you have no convergence rate. So neither family
gives you what you want.

The question of this talk is very simple: can we build a GAN-based solver
that handles exactly this Hamiltonian, and gives an explicit, sharp
convergence rate?

The answer is yes, in a precise regime that I will define on the next slide.

\section*{[Slide 6 -- The key insight]}

Here is the key insight. It is very simple, but it changes everything.

The observation is the following: \emph{if rho is held fixed, any non-separable
Hamiltonian becomes separable.}

Look at the diagram. On the left, we have the full non-separable Mean Field
Game with Hamiltonian $H$ of $x$, rho, $p$. We freeze rho at the current
iterate, called rho-$k$. Now the effective Hamiltonian, $H^k$ of $x$ and $p$,
is just $H_0$ plus epsilon times $R$ evaluated at rho-$k$. This is now a
function of only $x$ and $p$. It is separable in the sense the GAN solvers
need.

So we can apply APAC-Net to this separable subproblem. We get a new density,
rho-$k$-plus-one. Then we update rho-$k$ to rho-$k$-plus-one and we repeat.

This is essentially a fixed-point iteration on the population density. Each
step is one APAC-Net call.

The whole question of the rest of the talk is: when does this outer iteration
converge, and how fast?

\section*{[Slide 7 -- Weakly non-separable Hamiltonians]}

Now I want to be precise about the class of problems we can handle.

We say $H$ is \emph{weakly non-separable with strength epsilon} if it
decomposes into three parts.

First, a base Hamiltonian $H_0$. This part must be strongly convex in $p$,
with constant $c_0$.

Second, a Lasry-Lions monotone interaction $f_0$ of rho. The constant lambda
measures how monotone it is. Larger lambda means stronger restoring force.

Third, a coupling residual $R$. This is the part that introduces the
non-separability. We require $R$ to be Lipschitz in rho, with constant $L_R$,
and Lipschitz in $p$, with constant $L_R$-$p$. The number epsilon is the
\emph{strength} of the coupling.

When epsilon is zero, you get the classical separable Mean Field Game. When
$R$ equals rho times $p$, you recover the traffic flow model. When $R$ depends
only on rho, the constant $L_R$-$p$ is zero.

Now look at the box at the bottom. The most important quantity is this:
epsilon times $L_R$, strictly less than lambda.

We call this the \emph{GAN-tractable regime}. Intuitively, the coupling
strength must be controlled by the restoring force from the monotone
interaction. Below this threshold, our method works. Above it, it fails.
We will see both halves of this in the next slides.

\section*{[Slide 8 -- The Extended-MFGAN algorithm]}

The algorithm is essentially what I described before, written in pseudocode.

We initialise rho-zero to be the initial density. Then we loop.

In each iteration, we build the effective Hamiltonian $H^k$ by plugging the
current density rho-$k$ into the residual $R$. Then we run APAC-Net on the
resulting separable Mean Field Game. APAC-Net gives back a new density
rho-$k$-plus-one and a new value function. We update and repeat.

After $K$ outer iterations, we return the final density and value function.

The total cost is $K$ times one APAC-Net run.

How big does $K$ need to be? Once we prove our geometric convergence rate
kappa, the answer is given by a simple formula. To reach a tolerance delta,
you need $K$ equal to the ceiling of log of one over delta, divided by log
of one over kappa. So $K$ is small. It does not depend on the dimension,
only on how much accuracy you want.

A nice sanity check: when epsilon is zero, $H^k$ does not depend on rho-$k$
at all. The outer loop converges in one step. So Extended-MFGAN \emph{reduces
exactly to APAC-Net} in the separable limit. Our method is a true extension,
not a replacement.

In practice, the inner APAC-Net solve is the expensive part. Each call
trains a generator and a critic by adversarial training. So we want the
outer loop to converge in as few iterations as possible. The contraction
rate kappa controls exactly that. If kappa is small, very few outer
iterations are enough. If kappa is close to one, we need more iterations.
And if kappa is bigger than one, the iteration does not converge at all.
So the rate kappa is the central object in our analysis.

\section*{[Slide 9 -- The main result]}

Now to the main theorem.

Under the GAN-tractable condition, epsilon times $L_R$ less than lambda,
the Extended-MFGAN outer iteration is a contraction in $L^2$. The contraction
rate is exactly kappa equals epsilon times $L_R$, divided by lambda.

This means: the error after iteration $k$ plus one is at most kappa times
the error at iteration $k$. So the error decays geometrically, like kappa
to the $k$.

Because the densities are bounded, $L^2$ convergence implies Wasserstein-2
convergence at the same rate. So we have a strict improvement over MFDGM,
which only gets $L^p$ for $p$ less than two.

This is the first \emph{explicit} convergence rate for any deep learning
Mean Field Game solver in the non-separable setting. APAC-Net does not give
a rate. MFDGM does not give a rate. We do.

The rate is also a closed-form function of the data: epsilon, $L_R$, and
lambda. You can read off how many iterations you need.

In the next slide, I will sketch how the proof works.

\section*{[Slide 10 -- Proof sketch: Lasry-Lions duality]}

The proof technique is called Lasry-Lions duality. It is the same technique
used in the classical theory to prove uniqueness of the Mean Field Game
equilibrium. We adapt it to our perturbed non-separable setting.

The high-level idea is the following. We compare two outputs of the inner
solver, when we feed in two different frozen densities. We want to bound
the difference of the outputs by the difference of the inputs, with a small
constant. That constant will be our kappa.

The proof has six steps.

Take two Mean Field Game solutions, with densities rho-one and rho-two,
and value functions phi-one and phi-two. We compute the time derivative
of the inner product of delta-phi with delta-rho. Here delta means the
difference between the two solutions.

Step one: substitute the HJB equation and the Fokker-Planck equation for
both solutions.

Step two: the diffusion terms, nu times Laplacian, cancel exactly between
the two equations. This is the magic of the Lasry-Lions trick.

Step three: integrate by parts on the current term.

Step four: integrate from time zero to time $T$. The boundary terms vanish
because of the boundary conditions: delta-rho is zero at time zero, and
delta-phi is zero at time $T$.

Step five: split the resulting Hamiltonian-difference into two parts.
The first part comes from the convex base $H_0$. By strong convexity,
this part gives a non-positive contribution. The second part comes from
the coupling residual $R$. This is controlled by the Lipschitz constants.

Step six: combine with the Lasry-Lions monotonicity of $f_0$.

The final estimate is the boxed inequality at the bottom of the slide.
Lambda times the squared $L^2$ norm of delta-rho is at most epsilon times
$L_R$, times the product of the input and output $L^2$ norms.

Dividing by the $L^2$ norm of delta-rho, we get the contraction. Then by
the Banach Fixed-Point Theorem, the iteration converges geometrically with
rate kappa.

\section*{[Slide 11 -- Sharpness: rate is tight]}

So far, kappa is just an upper bound on the contraction rate. But actually
it is the exact rate. I want to convince you of this.

We prove this by a direct linearisation. Take the bilinear case, where $R$
equals rho times $p$. Take the ergodic Mean Field Game on the torus. The
unique uniform equilibrium has rho-star identically equal to one.

Now we linearise the outer fixed-point map at this point.

The stationary Fokker-Planck equation reduces to: Laplacian of delta-phi
equals zero. On the torus, this forces the gradient of delta-phi to vanish.

With this, the linearised HJB collapses to a simple scalar relation:
minus lambda times delta-rho, plus epsilon times $L_R$ times delta-rho-bar,
equals zero.

Solving: delta-rho equals epsilon $L_R$ over lambda, times delta-rho-bar.

So the operator norm of the linearised map is exactly kappa. The upper
bound in the theorem is attained. It is not just a bound. It is the
rate.

\section*{[Slide 12 -- Necessity: condition is required]}

The previous slide shows the rate is tight. This slide goes one step further:
the GAN-tractable condition is \emph{necessary}.

We use the same example. The same ergodic Mean Field Game. The same
linearisation. We just look at what happens when epsilon times $L_R$ is
larger than or equal to lambda.

If epsilon $L_R$ equals lambda, the operator norm of the linearised map is
exactly one. So the error stays the same forever. No convergence.

If epsilon $L_R$ is greater than lambda, the operator norm is greater than
one. The error grows geometrically. The iteration diverges. So the
condition is not just sufficient, it is also necessary.

This matters for practice. If you have a problem with epsilon $L_R$ above
the threshold, our method will not work. You should use MFDGM instead.
The boundary between the two regimes is exactly the GAN-tractable condition.

Look at the cartoon at the bottom. To the left of the threshold, you get
geometric decay. To the right, geometric growth. There is a sharp phase
transition at epsilon-star equals lambda over $L_R$.

So combining slide eleven and slide twelve: the GAN-tractable condition
is \emph{both necessary and sufficient} for linear convergence in the
linearised regime. This is unusual. Usually a contraction theorem only gives
you a bound. Ours gives you the actual rate.

\section*{[Slide 13 -- Experiment 1: separable limit]}

Now I want to show numerical results. Three experiments. Each one validates
a piece of the theory.

The first experiment is the simplest. The theory says that at epsilon equals
zero, one outer iteration is enough.

Look at the highlighted row in the table. At epsilon equals zero, the error
after one iteration is exactly zero. Down to machine precision. After eight
iterations, still zero.

For epsilon greater than zero, the error decays geometrically. The rate is
exactly equal to epsilon, just as the theory predicts.

This is a sanity check, but an important one. It confirms what the algorithm
description claimed: at epsilon equals zero, Extended-MFGAN reduces exactly
to one APAC-Net call.

\section*{[Slide 14 -- Experiment 2: bilinear MFG-LWR]}

The second experiment is the bilinear traffic flow model. Here the structure
is more interesting.

For $R$ equals rho times $p$, when you linearise at the uniform equilibrium,
the optimal momentum is one minus epsilon. So the Lipschitz constant $L_R$,
which is the absolute value of the momentum, is the absolute value of one
minus epsilon.

So the contraction rate is epsilon times the absolute value of one minus
epsilon, divided by lambda.

This formula is non-monotonic. It grows from zero, reaches a maximum at
epsilon equals one half, and goes back to zero at epsilon equals one.

For lambda equal to zero point two, the equation kappa equals one has two
roots in the interval zero to one: at about zero point two seven six, and
at about zero point seven two four.

You see this on the left panel. The empirical rate, shown as red dots,
sits exactly on the theoretical curve, shown as a black dashed line. Two
phase transitions, just as predicted.

On the right panel: the convergence curves for each value of epsilon.
Below the first threshold and above the second, you see geometric decay.
Between them, geometric growth.

Empirical and theoretical rates agree to four decimal places, across all
eleven values of epsilon.

I want to emphasise that this is not just a qualitative match. It is a
quantitative match to four digits. This is a strong validation of the
theory, and it is a feature of having an explicit rate. With a method
that gives you no rate, you cannot do this kind of comparison at all.

\section*{[Slide 15 -- Experiment 3: density-only phase transition]}

The third experiment is the cleanest validation. We use density-only
coupling, with both lambda and $L_R$ equal to one. So the predicted
threshold is at epsilon equals one.

We sweep epsilon from zero point one to one point five.

On the left panel, the blue dots are the empirical rate, and the black
dashed line is the theoretical kappa. They lie on top of each other.
The transition happens exactly at epsilon equals one, marked by the red
vertical line.

On the right panel, the convergence curves. For epsilon equals zero point
one, in the lightest colour, the error drops fourteen orders of magnitude
in twelve iterations. For epsilon equals one point five, in the brightest
colour, the error grows about an order of magnitude. The transition between
the two regimes is sharp.

Look at the small table at the bottom. At epsilon equals zero point one,
the final error is three point five times ten to the minus fourteen.
At epsilon equals one, the iteration stagnates. At epsilon equals one
point five, the error is four point six. Empirical kappa equals
theoretical kappa to all four digits.

\section*{[Slide 16 -- Where Extended-MFGAN sits]}

Let me put our method in context with prior work.

Top row, APAC-Net and MFGAN: they cannot handle non-separable Hamiltonians.
This is the red cell. They have distributional representation, Wasserstein-1
convergence, but the rate is implicit.

Middle row, MFDGM: handles any non-separable Hamiltonian, the green cell.
But pointwise neural network representation, only $L^p$ convergence for
$p$ less than two, and no rate. Three red cells.

Bottom row, Extended-MFGAN, our method: weakly non-separable Hamiltonian,
distributional rho through the generator, Wasserstein-2 convergence, and
explicit rate. All four green cells. We occupy the gap that no prior
method filled.

The trade-off is honest. We are restricted to the GAN-tractable regime.
Above the threshold, MFDGM remains the right tool. We do not replace
MFDGM, we complement it.

\section*{[Slide 17 -- Contributions in one slide]}

Let me summarise our contributions.

First: a new class of Hamiltonians. Weakly non-separable. It interpolates
between the separable and non-separable cases.

Second: a new algorithm. Extended-MFGAN. It decomposes a non-separable
Mean Field Game into a sequence of separable subproblems, each solved
by APAC-Net.

Third: a Wasserstein-2 contraction theorem. With explicit rate kappa
equals epsilon $L_R$ over lambda. The first explicit rate for any deep
Mean Field Game solver in this setting.

Fourth: that rate is sharp, and the condition is necessary. So our
theorem cannot be improved.

Fifth: numerical experiments where the empirical rate matches the theory
to four decimal places. Across three different benchmarks.

\section*{[Slide 18 -- Future directions]}

Briefly, what comes next.

First, the present numerics use a clean spectral scheme. A full neural
network implementation of the inner GAN solver is engineering work, but
a natural next step.

Second, with that implementation, a full pipeline-level benchmark against
MFDGM.

Third, scaling to high dimensions.

Fourth, a hybrid algorithm. Use the GAN inner solver when the GAN-tractable
condition holds. Fall back to MFDGM when it does not.

Fifth, extending the framework to mean field control problems, where a
social planner optimises over the whole population.

And sixth, applications to multi-agent reinforcement learning and to
engineering systems modelled by congestion-type Mean Field Games.

I think the most exciting direction is the hybrid algorithm. The two
families, GAN-based and physics-informed, do not need to compete. They
can work together. Below the threshold, use the GAN solver and enjoy the
explicit rate. Above the threshold, fall back to MFDGM. The threshold
itself is computable from the problem data, so the user can choose
automatically.

\section*{[Slide 19 -- Thank you]}

This brings me to the end of the talk.

To summarise in one sentence: we have shown how to combine the strengths
of GAN-based and physics-informed deep solvers for Mean Field Games, in
a precise weakly non-separable regime, with the first explicit and sharp
Wasserstein convergence rate.

The full paper, the code for all the experiments, and these slides are
available on GitHub. The link is on the screen. You can also reach me
by email.

Thank you for your attention. I am happy to take any questions.

\end{document}
